\phantomsection
\addcontentsline{toc}{section}{ANEXOS}

\newcounter{anexo}
\renewcommand{\theanexo}{\Alph{anexo}}
\newcommand{\anexotitulo}[2]{%
  \refstepcounter{anexo}%
  \setcounter{table}{0}%
  \renewcommand{\thetable}{\theanexo.\arabic{table}}%
  \subsection*{Anexo~\theanexo. #1}%
  \phantomsection
  \addcontentsline{toc}{subsection}{Anexo~\theanexo. #1}%
  \label{#2}%
}

\anexotitulo{Diccionario de atributos territoriales}{anexo:diccionario_atributos}

Este anexo documenta los 19 atributos de localización que integran el vector de entrada aprobado para el modelo. Las variables diagnósticas se reportan por separado y el flujo observado constituye la variable objetivo. Las transformaciones se ajustan únicamente con los orígenes activos del conjunto de entrenamiento de cada rama y escenario.

\begin{table}[H]
  \centering
  \scriptsize
  \caption{Atributos de población, uso de suelo y red vial del vector de entrada.}
  \label{tab:anexo_atributos_base}
  \begin{tabularx}{\textwidth}{@{}p{0.34\textwidth}p{0.42\textwidth}X@{}}
    \toprule
    \textbf{Atributo} & \textbf{Fuente y construcción} & \textbf{Unidad y transformación} \\
    \midrule
    \texttt{population\_census\_2024\_per\_km2} &
    Censo 2024 del INE; asignación areal de población y división por el área métrica de la unidad. &
    personas/km$^2$; $\log(1+x)$ y estandarización $z$. \\
    \texttt{landuse\_residential\_share} &
    OpenStreetMap (corte 2024-01-01); unión de polígonos residenciales recortados y división por el área de la unidad. &
    km$^2$/km$^2$; estandarización $z$. \\
    \texttt{landuse\_commercial\_share} &
    OpenStreetMap (corte 2024-01-01); unión de polígonos comerciales recortados y división por el área de la unidad. &
    km$^2$/km$^2$; estandarización $z$. \\
    \texttt{landuse\_industrial\_share} &
    OpenStreetMap (corte 2024-01-01); unión de polígonos industriales recortados y división por el área de la unidad. &
    km$^2$/km$^2$; estandarización $z$. \\
    \texttt{landuse\_retail\_share} &
    OpenStreetMap (corte 2024-01-01); unión de polígonos de comercio minorista recortados y división por el área de la unidad. &
    km$^2$/km$^2$; estandarización $z$. \\
    \texttt{landuse\_natural\_share} &
    OpenStreetMap (corte 2024-01-01); unión de polígonos naturales recortados y división por el área de la unidad. &
    km$^2$/km$^2$; estandarización $z$. \\
    \texttt{road\_residential\_km\_per\_km2} &
    OpenStreetMap; unión de líneas residenciales recortadas y división por el área de la unidad. &
    km/km$^2$; $\log(1+x)$ y estandarización $z$. \\
    \texttt{road\_main\_km\_per\_km2} &
    OpenStreetMap; unión de líneas principales recortadas y división por el área de la unidad. &
    km/km$^2$; $\log(1+x)$ y estandarización $z$. \\
    \texttt{road\_other\_km\_per\_km2} &
    OpenStreetMap; unión de las demás líneas viales recortadas y división por el área de la unidad. &
    km/km$^2$; $\log(1+x)$ y estandarización $z$. \\
    \bottomrule
  \end{tabularx}
\end{table}

\begin{table}[H]
  \centering
  \scriptsize
  \caption{Atributos de puntos de interés y superficie por familia temática.}
  \label{tab:anexo_atributos_poi}
  \begin{tabularx}{\textwidth}{@{}p{0.31\textwidth}p{0.43\textwidth}X@{}}
    \toprule
    \textbf{Atributo} & \textbf{Fuente y construcción} & \textbf{Unidad y transformación} \\
    \midrule
    \texttt{transport\_poi\_per\_km2} &
    OpenStreetMap; nodos estrictamente interiores, deduplicados por homónimos nodo--área, categoría y coincidencia espacial, divididos por el área de la unidad. &
    objetos/km$^2$; $\log(1+x)$ y estandarización $z$. \\
    \texttt{transport\_area\_share} &
    OpenStreetMap; unión de las superficies de transporte recortadas y división por el área de la unidad. &
    km$^2$/km$^2$; estandarización $z$. \\
    \texttt{food\_poi\_per\_km2} &
    OpenStreetMap; nodos estrictamente interiores, deduplicados por homónimos nodo--área, categoría y coincidencia espacial, divididos por el área de la unidad. &
    objetos/km$^2$; $\log(1+x)$ y estandarización $z$. \\
    \texttt{food\_area\_share} &
    OpenStreetMap; unión de las superficies de alimentación recortadas y división por el área de la unidad. &
    km$^2$/km$^2$; estandarización $z$. \\
    \texttt{health\_poi\_per\_km2} &
    OpenStreetMap; nodos estrictamente interiores, deduplicados por homónimos nodo--área, categoría y coincidencia espacial, divididos por el área de la unidad. &
    objetos/km$^2$; $\log(1+x)$ y estandarización $z$. \\
    \texttt{health\_area\_share} &
    OpenStreetMap; unión de las superficies de salud recortadas y división por el área de la unidad. &
    km$^2$/km$^2$; estandarización $z$. \\
    \texttt{education\_poi\_per\_km2} &
    OpenStreetMap; nodos estrictamente interiores, deduplicados por homónimos nodo--área, categoría y coincidencia espacial, divididos por el área de la unidad. &
    objetos/km$^2$; $\log(1+x)$ y estandarización $z$. \\
    \texttt{education\_area\_share} &
    OpenStreetMap; unión de las superficies educacionales recortadas y división por el área de la unidad. &
    km$^2$/km$^2$; estandarización $z$. \\
    \texttt{retail\_poi\_per\_km2} &
    OpenStreetMap; nodos estrictamente interiores, deduplicados por homónimos nodo--área, categoría y coincidencia espacial, divididos por el área de la unidad. &
    objetos/km$^2$; $\log(1+x)$ y estandarización $z$. \\
    \texttt{retail\_area\_share} &
    OpenStreetMap; unión de las superficies de comercio minorista recortadas y división por el área de la unidad. &
    km$^2$/km$^2$; estandarización $z$. \\
    \bottomrule
  \end{tabularx}
\end{table}

La distancia $d_{ij}$ se calcula como distancia geodésica esférica WGS84 entre los centroides representativos de los pares origen--destino y se incorpora al vector como medida de separación espacial. Una ruta por la red requiere información adicional de infraestructura y operación. Las capas territoriales pueden superponerse, por lo que sus proporciones se analizan por familia. Para OpenStreetMap, un valor cero registra ausencia de objetos mapeados conforme a la regla aplicada.

\clearpage
\anexotitulo{Trazabilidad de las corridas mensuales}{anexo:corridas_mensuales}

Este anexo individualiza las nueve corridas del experimento mensual: tres inicializaciones independientes en cada representación. La época seleccionada corresponde al mayor CPC global de validación de la trayectoria, conservando la primera época ante empate; la época final indica cuándo se detuvo la corrida. Las dos continuaciones autorizadas partieron del límite inicial de 200 épocas y conservaron datos, particiones, atributos, arquitectura, optimizador y regla de selección. Todas las corridas cerraron por meseta de validación.

\begin{table}[H]
  \centering
  \scriptsize
  \caption{Corridas del experimento mensual y selección por validación.}
  \label{tab:anexo_corridas}
  \resizebox{\textwidth}{!}{%
    \begin{tabular}{@{}lrrrrll@{}}
      \toprule
      \textbf{Representación} & \textbf{Semilla} & \textbf{Época seleccionada} & \textbf{Época final} & \textbf{CPC validación} & \textbf{Extensión} & \textbf{Cierre} \\
      \midrule
      Zona 777 & 20260909 & 196 & 226 & 0,532453 & 200--226 & Meseta \\
      Zona 777 & 20260910 & 94  & 124 & 0,525924 & ---      & Meseta \\
      Zona 777 & 20260911 & 74  & 104 & 0,517290 & ---      & Meseta \\
      H3-r7    & 20260909 & 72  & 97  & 0,737286 & ---      & Meseta \\
      H3-r7    & 20260910 & 162 & 192 & 0,734555 & ---      & Meseta \\
      H3-r7    & 20260911 & 69  & 99  & 0,732907 & ---      & Meseta \\
      H3-r8    & 20260909 & 158 & 188 & 0,600163 & ---      & Meseta \\
      H3-r8    & 20260910 & 188 & 218 & 0,594110 & 200--218 & Meseta \\
      H3-r8    & 20260911 & 137 & 167 & 0,589980 & ---      & Meseta \\
      \bottomrule
    \end{tabular}%
  }
\end{table}

La prueba se evaluó después de fijar y congelar los \textit{checkpoints} seleccionados por validación. Los resultados de las tres semillas se presentan completos en el Capítulo 5; esta tabla conserva la trazabilidad de la selección.

\clearpage
\anexotitulo{Controles de soporte, conservación de masa y ejecución}{anexo:controles_ejecucion}

Este anexo reúne los controles de reproducibilidad y el detalle de los recursos medidos junto con los denominadores de la repetición independiente, la conservación del flujo por origen, el detalle por corrida y las incidencias técnicas resueltas.

\subsubsection*{Soporte de prueba, repetición independiente y conservación de masa}

\begin{table}[H]
  \centering
  \scriptsize
  \caption{Controles de soporte, repetición independiente y conservación de masa.}
  \label{tab:anexo_controles_soporte}
  \resizebox{\textwidth}{!}{%
    \begin{tabular}{@{}lrrrrrr@{}}
      \toprule
      \textbf{Representación} & \textbf{Orígenes/destinos prueba} & \textbf{Unidades sin origen} & \textbf{CPC congelado} & \textbf{CPC repetido} & \textbf{Diferencia CPC} & \textbf{Máx. error por origen} \\
      \midrule
      Zona 777 & 101/793  & 0  & 0,5647245226 & 0,5647245225 & $9{,}576\times10^{-11}$ & $9{,}313\times10^{-10}$ \\
      H3-r7    & 32/223   & 5  & 0,7670518160 & 0,7670518160 & $3{,}483\times10^{-11}$ & $1{,}863\times10^{-9}$ \\
      H3-r8    & 189/1210 & 21 & 0,6132965073 & 0,6132965073 & $9{,}350\times10^{-12}$ & $1{,}630\times10^{-9}$ \\
      \bottomrule
    \end{tabular}%
  }
\end{table}

\noindent\footnotesize\textit{Nota.} El error máximo corresponde al máximo entre las tres semillas y los dos esquemas de ponderación. La repetición independiente usa la semilla 20260909 y el conjunto completo de destinos. El soporte de prueba comprende los orígenes asignados al bloque de prueba; las unidades sin origen se registran por separado y se excluyen del cálculo del CPC.\normalsize

La repetición independiente recuperó el CPC congelado de cada representación, con diferencias inferiores a $10^{-6}$. Las nueve corridas preservaron el flujo por origen; el mayor error observado fue $1{,}863\times10^{-9}$. Estos controles verifican la consistencia técnica de la ejecución; la evaluación espacial y temporal se desarrolla mediante el protocolo del Capítulo 5.

\subsubsection*{Recursos medidos por corrida}

\begin{table}[H]
  \centering
  \small
  \caption{Tiempo de entrenamiento, memoria RAM y tamaño de artefactos por corrida.}
  \label{tab:anexo_recursos}
  \begin{tabular}{@{}lrrrr@{}}
    \toprule
    \textbf{Representación} & \textbf{Semilla} & \textbf{Entrenamiento (h)} & \textbf{RAM máx. (GiB)} & \textbf{Artefactos (GiB)} \\
    \midrule
    Zona 777 & 20260909 & 1,30 & 0,53 & 1,16 \\
    Zona 777 & 20260910 & 0,72 & 0,53 & 0,64 \\
    Zona 777 & 20260911 & 0,60 & 0,53 & 0,53 \\
    H3-r7    & 20260909 & 0,05 & 0,53 & 0,50 \\
    H3-r7    & 20260910 & 0,10 & 0,53 & 0,99 \\
    H3-r7    & 20260911 & 0,05 & 0,53 & 0,51 \\
    H3-r8    & 20260909 & 2,19 & 0,78 & 0,97 \\
    H3-r8    & 20260910 & 2,50 & 0,83 & 1,12 \\
    H3-r8    & 20260911 & 1,95 & 0,78 & 0,86 \\
    \midrule
    Suma/máximo (9 corridas) & --- & 9,46 & 0,83 & 7,27 \\
    \bottomrule
  \end{tabular}
\end{table}

Las horas corresponden a las épocas de entrenamiento y validación. La carga inicial y la escritura de \textit{checkpoints} se registran fuera de esa medida, por lo que las horas son una cota inferior del tiempo total de ejecución. La RAM se informa como máximo observado y el tamaño corresponde a los artefactos conservados de cada corrida. Las condiciones de ejecución difieren entre ramas, por lo que la tabla reporta tiempo de entrenamiento y validación en vez de tiempo calendario.

\subsubsection*{Incidencias técnicas resueltas}

\begin{table}[H]
  \centering
  \small
  \caption{Incidencias técnicas y efecto sobre los resultados reportados.}
  \label{tab:anexo_incidencias}
  \begin{tabularx}{\textwidth}{@{}p{0.24\textwidth}p{0.48\textwidth}X@{}}
    \toprule
    \textbf{Alcance} & \textbf{Incidencia y resolución} & \textbf{Efecto} \\
    \midrule
    Zona 777, semilla 20260909 &
    Se detectó la ausencia del parámetro de hilos de CPU antes de iniciar el entrenamiento. Se corrigió antes de entrenar y evaluar la corrida. &
    La corrección se realizó antes del entrenamiento y la evaluación; la corrida conservó su plan y sus resultados. \\
    Cierre y repetición independiente &
    Una ruta relativa de serialización se detectó luego de generar productos parciales. La repetición se ejecutó en un directorio nuevo y se excluyeron los productos parciales. &
    Los \textit{checkpoints} congelados y la evaluación de prueba conservaron los resultados reportados. \\
    \bottomrule
  \end{tabularx}
\end{table}
